\section{Conclusions}
\label{sec:conclusions}

Procurement transport planning with returnable transport items has
been formulated with transport capacity as what it physically is --- a
three-dimensional loading condition on every dispatched truck --- and
solved by branch and price over complete loaded trucks. Because
3D-SCPTOP is written over trucks it is block-angular, so the
Dantzig--Wolfe master is obtained by aggregation and carries no loading
configuration (Table~\ref{tab:rowmap}), and the equivalence needs no
hypothesis. Every lower bound reported rests on
Proposition~\ref{prop:inclusion}, which makes a bound affordable where
a geometric certificate is not. And a volumetric transport capacity is
not a conservative simplification: the method plans the same
quantities with $27.0\%$ and $24.5\%$ fewer trucks than re-packing two
of its plans needs, and on the industrial catalogue the error runs one
way only.

Three limits are real, and each names its next step. The method
certifies no optimum here: the search ends with regions unrefuted, and the
one optimum reported is established by inclusion from an exactly solved
relaxation. The bound is loose exactly where the crush limit halves a
container family's layer count, and the stack-count rows
\eqref{eq:stackrows} tighten it measurably, so they are the candidate
for tightening it inside the search as well, once it is known whether
the loads they admit can be placed. Closing the regions the search
leaves unrefuted may need a branching rule on \emph{pairs} of columns,
untested here. And the comparison against a procedure in industrial use is to
be run on instances a plant solves itself, against a plan supplied to
it.


